\documentclass[11pt]{article}

\usepackage[a4paper,margin=28mm]{geometry}
\usepackage{amsmath,amssymb,amsthm,mathtools}
\usepackage[T1]{fontenc}
\usepackage{lmodern}
\usepackage{microtype}
\usepackage[hidelinks]{hyperref}
\usepackage{enumitem}

\newcommand{\NE}{\overline{\operatorname{NE}}}
\newcommand{\Nef}{\operatorname{Nef}}
\newcommand{\Mov}{\operatorname{Mov}}
\newcommand{\Proj}{\operatorname{Proj}}
\newcommand{\Spec}{\operatorname{Spec}}
\newcommand{\Pic}{\operatorname{Pic}}
\newcommand{\Cl}{\operatorname{Cl}}
\newcommand{\Exc}{\operatorname{Exc}}
\newcommand{\cO}{\mathcal O}

\title{The Absolute Two-Ray Game as a Computational Testbed for the MMP}
\author{}
\date{July 2026}

\begin{document}
\maketitle

\begin{abstract}
This note records a discussion about using the absolute two-ray game as a
computational testbed for a general algorithm for the Minimal Model Program
(MMP).  The intended method does not use Cox rings or variation of GIT to
solve the problem.  Instead, it uses operations that plausibly extend to a
general MMP algorithm: nefness testing, extremal-ray detection, construction
of contractions by section rings, construction of flips by relative canonical
algebras, and saturation in Rees-algebra or graph calculations.  A terminal
quartic threefold with a $cA_5$ point is proposed as a concrete example with
two flips.
\end{abstract}

\section{Absolute Picard rank two}

Let $X$ be a projective variety over a field of characteristic zero.  In the
absolute setting the base is $\Spec k$.  If
\[
  \rho(X)=\dim_{\mathbb R}N_1(X)_{\mathbb R}=2,
\]
then $\NE(X)$ is a closed convex cone in a two-dimensional real vector space.
An ample divisor is strictly positive on nonzero effective curve classes, so
$\NE(X)$ is pointed.  Consequently it has two boundary rays:
\[
  \NE(X)=R_1+R_2.
\]
Thus Picard rank two really does reduce the numerical geometry to two rays.

There is nevertheless an important distinction between the following
statements:
\begin{enumerate}[label=(\roman*)]
  \item the two boundary directions exist as real rays;
  \item a boundary ray is rational;
  \item it is generated by the class of an actual curve;
  \item it is $(K_X+\Delta)$-negative and admits an extremal contraction.
\end{enumerate}
For a general projective variety, a boundary ray can be irrational and need
not contain a nonzero integral curve class.  Under the hypotheses of the cone
and contraction theorems, a $(K_X+\Delta)$-negative extremal ray is generated
by a rational curve and admits a contraction.

The dual cone is
\[
  \Nef(X)=\NE(X)^\vee\subset N^1(X)_{\mathbb R}.
\]
Therefore computing the Mori cone and computing the nef cone are dual versions
of the same problem.  Picard rank two reduces the dimension, but it does not
by itself reveal the slopes of the two boundary rays.

\section{The Fano simplification}

Assume that $X$ is $\mathbb Q$-factorial and klt, that $\rho(X)=2$, and that
$-K_X$ is ample.  Both boundary rays of $\NE(X)$ are then $K_X$-negative.
The cone theorem gives
\[
  \NE(X)=\mathbb R_{\geq 0}[C_1]
          +\mathbb R_{\geq 0}[C_2],
\]
where $C_1$ and $C_2$ may be chosen to be rational curves.  One also has a
length bound of the form
\[
  0<-K_X\cdot C_i\leq 2\dim X.
\]
If $mK_X$ is Cartier and $H=-mK_X$, then
\[
  1\leq H\cdot C_i\leq 2m\dim X.
\]
This turns the theoretical search for the two rays into a finite bounded-degree
search in an anticanonical embedding.  Such a Hilbert- or Chow-scheme search is
still potentially expensive, but irrational boundary rays no longer occur.

Once $[C_1]$ and $[C_2]$ are known, nefness of a divisor $D$ is reduced to
two inequalities:
\[
  D\text{ is nef}
  \quad\Longleftrightarrow\quad
  D\cdot C_1\geq 0\ \text{ and }\ D\cdot C_2\geq 0.
\]
For a positive-dimensional Fano variety, $K_X$ itself is of course not nef.
The nontrivial nefness tests arise on extracted or flipped intermediate models,
or for divisors other than $K_X$.

\section{Recovering a ray from a nef threshold}

Choose an ample divisor $A$ and a direction $D\in N^1(X)_{\mathbb R}$.  Define
\[
  \tau=\sup\{t\geq 0\mid A+tD\text{ is nef}\}.
\]
Then
\[
  L=A+\tau D
\]
lies on the boundary of the nef cone.  Since $\rho(X)=2$, the orthogonal space
\[
  L^\perp\subset N_1(X)_{\mathbb R}
\]
is one-dimensional and gives the corresponding boundary ray of $\NE(X)$.

If a rationality theorem applies and an explicit denominator bound $q\leq Q$
is available for $\tau=p/q$, a nefness decision procedure can be used as a
monotone oracle.  Binary search to precision less than $1/(2Q^2)$, followed by
rational reconstruction, determines $\tau$ exactly.  The curve direction is
then obtained by solving the single linear equation
\[
  L\cdot\gamma=0.
\]
This is one possible way in which Picard rank two can make an existing general
nefness algorithm substantially more useful.

The numerical class $\gamma$ is not yet an actual embedded curve.  When $L$ is
semiample, one can construct the morphism defined by $|rL|$ for sufficiently
divisible $r$ and obtain curves from its positive-dimensional fibers.

\section{A Cox-free computational MMP}

The aim of the proposed experiment is not to compute the chamber decomposition
of a Cox ring.  Cox rings may be used externally to check an answer, but they
should not be part of the algorithm under test.  A generalizable step should
instead have the following form.

Let $Y_i$ be a $\mathbb Q$-factorial klt model of Picard rank two.
\begin{enumerate}
  \item Use nefness testing and a one-dimensional threshold computation to
        identify a $(K_{Y_i}+\Delta_i)$-negative extremal ray $R_i$.
  \item Find a supporting nef divisor $L_i$ satisfying
        \[
          L_i\cdot R_i=0.
        \]
  \item Construct the contraction from a section ring:
        \[
          g_i\colon Y_i\longrightarrow Z_i
          =\Proj\bigoplus_{m\geq 0}H^0(Y_i,mL_i).
        \]
  \item Compute the exceptional locus and decide whether $g_i$ is of fiber
        type, divisorial, or small.
  \item If $g_i$ is small, construct the flip from the relative canonical
        algebra.  If $rK_{Y_i}$ is Cartier, use
        \[
          Y_i^+
          =\Proj_{Z_i}
            \bigoplus_{m\geq 0}
            g_{i*}\cO_{Y_i}(mrK_{Y_i}).
        \]
  \item Compute a presentation of $Y_i^+$, update divisor and intersection
        data, and repeat.
\end{enumerate}

The difficult algorithmic questions are not merely the number of variables in
the equations.  They include certification that the relevant section or
canonical algebra has been generated, exact detection of the extremal ray, and
verification of the singularities after a flip.

\section{Elementary examples and why they are insufficient}

The blow-up of projective three-space at one point is a useful sanity check:
\[
  X=\operatorname{Bl}_p\mathbb P^3,
  \qquad p=[1:0:0:0].
\]
In $\mathbb P^3_x\times\mathbb P^2_u$ it is defined by
\[
  x_1u_2-x_2u_1=x_1u_3-x_3u_1=x_2u_3-x_3u_2=0.
\]
Writing $H$ for the pullback of a hyperplane and $E$ for the exceptional
divisor, one has
\[
  -K_X=4H-2E=2H+2(H-E),
\]
and
\[
  \Nef(X)=\mathbb R_{\geq0}[H]
           +\mathbb R_{\geq0}[H-E].
\]
The two contractions are the blow-down $X\to\mathbb P^3$ and the
$\mathbb P^1$-bundle $X\to\mathbb P^2$.  Thus either choice terminates in one
step; no flip occurs.

This is typical of simple smooth Fano threefold examples.  To obtain flips in
dimension three, it is natural to allow terminal singularities and to start
with a divisorial extraction from a Fano threefold of Picard rank one.  The
extracted model has Picard rank two and is the actual input to the absolute
two-ray game.

\section{A proposed quartic threefold test case}

Consider the quartic threefold
\[
  X=X_{2,2}\subset\mathbb P^4_{x_0,\ldots,x_4}
\]
defined by
\begin{align*}
0={}&(x_0x_1-x_3^2)^2-(x_0x_2-x_4^2)^2\\
   &+x_1^2x_3^2+x_2^2x_4^2+x_1^4+x_2^4.
\end{align*}
It is a terminal factorial Fano threefold whose only singular point is
\[
  P=[1:0:0:0:0],
\]
of type $cA_5$ \cite{AK}.

Set
\begin{align*}
  \alpha&=x_0(x_1+x_2)-(x_3^2+x_4^2),\\
  \beta &=x_0(x_1-x_2)-(x_3^2-x_4^2).
\end{align*}
Up to harmless rescaling of $\alpha$ and $\beta$, the equation is
\[
  \alpha\beta+x_1^2x_3^2+x_2^2x_4^2+x_1^4+x_2^4=0,
\]
and the germ at $P$ has the form
\[
  \alpha\beta+x_3^6+x_4^6+	ext{higher-weight terms}=0.
\]

Take the discrepancy-one weighted blow-up with local weights
\[
  \operatorname{wt}(\alpha,\beta,x_3,x_4)=(3,3,1,1).
\]
Denote the extracted variety by
\[
  f\colon Y_0\longrightarrow X.
\]
Then $\rho(Y_0)=2$.  The known Sarkisov link has two flips and terminates in a
degree-four del Pezzo fibration over $\mathbb P^1$ \cite{AK}.  The local flips
are hypersurface flips of type
\[
  (3,1,1,-1,-1;2).
\]
These facts should be treated as expected output for verification, not as
input to the algorithm.

The extraction has a small weighted Rees presentation.  With one exceptional
parameter $u$, a convenient strict-transform presentation is
\begin{align*}
  u\alpha
    &=x_0(x_1+x_2)-(x_3^2+x_4^2),\\
  u\beta
    &=x_0(x_1-x_2)-(x_3^2-x_4^2),\\
  \alpha\beta
    &+u^4(x_1^2x_3^2+x_2^2x_4^2)
     +u^2(x_1^4+x_2^4)=0.
\end{align*}
Thus the initial quartic has five variables and one equation, while the
extracted model can be represented using eight variables and three sparse
equations.  This is a plausible size for computation in systems such as
Macaulay2 or Singular.

\section{Expected computational behavior}

The equations need not grow exponentially if every step is reduced to a
minimal presentation.  A local flip of type $(3,1,1,-1,-1;2)$ is a
hypersurface calculation in five local variables.  A global relative-Proj
calculation may add one or two generators, but redundant old generators should
then be eliminated immediately.

The following implementation choices are likely to cause an explosion:
\begin{itemize}
  \item retaining the equations of every graph of every rational map;
  \item introducing new generators without eliminating old ones;
  \item replacing weighted or relative Proj presentations by large ordinary
        projective Veronese embeddings;
  \item computing section rings to a large degree in a single undifferentiated
        step;
  \item postponing all saturation and elimination until the end.
\end{itemize}

A better workflow is
\[
  \begin{aligned}
  \text{add minimal relative generators}
    &\longrightarrow \text{compute relations}\\
    &\longrightarrow \text{saturate}
     \longrightarrow \text{eliminate redundant variables}.
  \end{aligned}
\]
This is not a Cox-ring method; it is the ordinary algebraic implementation of
Rees algebras, graph closures, and relative Proj.

The test can be staged as follows.
\begin{description}
  \item[Stage 1.] Supply the extremal ray and supporting divisor, and test only
        the construction of the two flips.
  \item[Stage 2.] Supply the Picard and intersection data, but require the
        algorithm to discover the second ray using nefness tests.
  \item[Stage 3.] Supply only the extracted model $Y_0$ and require automatic
        ray detection, contraction, smallness detection, flip construction,
        and singularity verification.
\end{description}
Stage 1 should be comfortably feasible.  Stage 2 is also realistic.  Stage 3
is a meaningful stress test for a general MMP algorithm; the likely bottleneck
is certification of the relevant graded algebras and singularities rather than
the raw number of equations.

\section{Saturation by an exceptional or base ideal}

Saturation is used to remove components introduced artificially by pullback,
denominator clearing, or incidence equations.

Suppose a blow-up chart has exceptional divisor $E=\{u=0\}$.  If the total
transform of an equation is
\[
  u^m\widetilde f=0,
\]
then its zero scheme contains the whole exceptional divisor as an extra
component.  The strict-transform ideal is
\[
  I_{\mathrm{str}}
  =I_{\mathrm{tot}}:(u)^\infty
  =\bigcup_{n\geq0}\bigl(I_{\mathrm{tot}}:(u)^n\bigr).
\]
More generally,
\[
  I_{\mathrm{str}}
  =I_{\mathrm{tot}}:\mathcal I_E^\infty.
\]

For example, blow up the origin of $\mathbb A^2$ and use the chart
\[
  x=u,\qquad y=uv.
\]
The pullback of the curve $y-x^2=0$ is
\[
  u(v-u)=0.
\]
Saturation gives
\[
  (u(v-u)):(u)^\infty=(v-u),
\]
which is the strict transform.  Its intersection with $E$ is retained; only
the unwanted component $E$ itself is removed.

The same principle applies to graph constructions.  If a rational function is
represented by
\[
  s=\frac{f}{u^a},
\]
then clearing denominators gives $u^as-f=0$, which may acquire spurious
components on $u=0$.  Saturating by $(u)$ extracts the closure of the graph
over $u\neq0$.  More generally, for a rational map with base ideal $J$, one
often computes
\[
  I_{\Gamma}=I_{\mathrm{inc}}:J^\infty.
\]

The phrase ``saturate by the exceptional ideal'' is therefore shorthand.  One
should not blindly use the same ideal at every flip.  The correct ideal is
determined by the open locus whose closure is being computed: it may be an
exceptional-divisor ideal, a base ideal, or an irrelevant ideal in a chosen
relative-Proj presentation.

\section{Conclusion}

The absolute two-ray game is a promising intermediate test between toy
examples and a general MMP implementation.  Picard rank two removes branching
in the numerical cone, while terminal quartic examples retain genuinely hard
operations: extremal-ray detection, small contractions, flips, and singularity
control.  The $cA_5$ quartic above appears to have an appropriate balance: the
input equations are sparse and small, but the expected computation contains
two nontrivial flips.  Keeping relative-Proj and Rees presentations minimal is
essential for avoiding artificial equation growth.

\begin{thebibliography}{9}

\bibitem{AK}
H.~Ahmadinezhad and A.-S.~Kaloghiros,
\emph{Non-rigid quartic $3$-folds},
Compositio Math. \textbf{152} (2016), 955--983.
\url{https://arxiv.org/abs/1310.5554}

\bibitem{Campo}
L.~Campo,
\emph{Sarkisov links for index $1$ Fano $3$-folds in codimension $4$},
Math. Nachr. \textbf{296} (2023), 957--979.
\url{https://arxiv.org/abs/2011.12209}

\bibitem{Fujino}
O.~Fujino,
\emph{Fundamental theorems for the log minimal model program},
Publ. Res. Inst. Math. Sci. \textbf{47} (2011), 727--789.
\url{https://arxiv.org/abs/0909.4445}

\bibitem{HM}
C.~D.~Hacon and J.~McKernan,
\emph{The Sarkisov program},
J. Algebraic Geom. \textbf{22} (2013), 389--405.
\url{https://arxiv.org/abs/0905.0946}

\end{thebibliography}

\end{document}
